% Part: first-order-logic
% Chapter: sequent-calculus
% Section: quantifier-rules

\documentclass[../../../include/open-logic-section]{subfiles}

\begin{document}

\olfileid[id]{fol}{seq}{qrl}

\olsection{Aturan Kuantor}

\subsection{Aturan untuk $\lforall$}

\begin{defish}
\Axiom$ !A(t), \Gamma \fCenter \Delta$
\RightLabel{\LeftR{\lforall}}
\UnaryInf$ \lforall[x][!A(x)],\Gamma \fCenter \Delta$
\DisplayProof
\hfill
\Axiom$ \Gamma \fCenter \Delta, !A(a) $
\RightLabel{\RightR{\lforall}}
\UnaryInf$ \Gamma \fCenter \Delta, \lforall[x][!A(x)]$
\DisplayProof
\end{defish}

Dalam \LeftR{\lforall}, $t$ adalah suku tertutup (yakni suku tanpa
variabel). Dalam \RightR{\lforall}, $a$~adalah !!a{constant} yang tidak
boleh muncul di mana pun dalam sekuen bawah aturan \RightR{\lforall}.
Kita menyebut $a$ sebagai \emph{variabel eigen} dari inferensi
\RightR{\forall}.\footnote{Kita menggunakan istilah ``variabel eigen''
meskipun $a$ dalam aturan di atas adalah !!a{constant}. Alasannya bersifat
historis.}

\subsection{Aturan untuk $\lexists$}

\begin{defish}
\Axiom$ !A(a), \Gamma \fCenter \Delta $
\RightLabel{\LeftR{\lexists}}
\UnaryInf$ \lexists[x][!A(x)], \Gamma \fCenter \Delta$
\DisplayProof
\hfill
\Axiom$ \Gamma \fCenter \Delta, !A(t) $
\RightLabel{\RightR{\lexists}}
\UnaryInf$ \Gamma \fCenter \Delta, \lexists[x][!A(x)]$
\DisplayProof
\end{defish}

Sekali lagi, $t$~adalah suku tertutup, dan $a$~adalah !!a{constant}
yang tidak muncul dalam sekuen bawah aturan \LeftR{\lexists}. Kita
menyebut $a$ sebagai \emph{variabel eigen} dari inferensi
\LeftR{\lexists}.

Syarat bahwa suatu variabel eigen tidak muncul dalam sekuen bawah
inferensi \RightR{\lforall} atau \LeftR{\lexists} disebut
\emph{syarat variabel eigen}.

\begin{explain}
Ingat kembali konvensi bahwa ketika $!A$ adalah !!a{formula} dengan
!!{variable}~$x$ bebas, kita menandainya dengan menulis~$!A(x)$. Dalam
konteks yang sama, $!A(t)$ adalah singkatan dari~$\Subst{!A}{t}{x}$.
Jadi, aturan \RightR{\lexists} juga dapat kita tulis sebagai:
\begin{prooftree}
  \Axiom$\Gamma \fCenter \Delta, \Subst{!A}{t}{x}$
  \RightLabel{\RightR{\lexists}}
  \UnaryInf$\Gamma \fCenter \Delta, \lexists[x][!A]$
\end{prooftree}
Perhatikan bahwa $t$ mungkin sudah muncul dalam~$!A$; misalnya, $!A$
mungkin berupa~$\Atom{\Obj P}{t,x}$. Jadi, menyimpulkan $\Gamma \Sequent
\Delta, \lexists[x][\Atom{\Obj P}{t,x}]$ dari~$\Gamma \Sequent \Delta,
\Atom{\Obj P}{t,t}$ merupakan penerapan \RightR{\lexists} yang benar---kita
dapat ``mengganti'' satu atau lebih, dan tidak harus semua, kemunculan~$t$
dalam premis dengan !!{variable}~$x$ yang terikat. Akan tetapi, syarat
variabel eigen dalam \RightR{\lforall} dan~\LeftR{\lexists} mengharuskan
!!{constant}~$a$ tidak muncul dalam~$!A$. Karena itu, kita tidak dapat
menyimpulkan $\Gamma \Sequent \Delta, \lforall[x][\Atom{\Obj P}{a,x}]$
secara benar dari $\Gamma \Sequent \Delta, \Atom{\Obj P}{a,a}$ dengan
menggunakan~$\RightR{\lforall}$.
\end{explain}

\begin{explain}
Dalam \RightR{\lexists} dan \LeftR{\lforall} tidak ada pembatasan atas
suku~$t$. Sebaliknya, dalam aturan \LeftR{\lexists} dan
\RightR{\lforall}, syarat variabel eigen mengharuskan !!{constant}~$a$
tidak muncul di mana pun di luar~$!A(a)$ dalam sekuen atas. Syarat ini
diperlukan untuk memastikan sistem itu sahih, yakni hanya !!{derive}
sekuen-sekuen yang valid. Tanpa syarat ini, hal berikut akan diizinkan:
\begin{prooftree}
  \Axiom$!A(a) \fCenter !A(a)$
  \RightLabel{*\LeftR{\lexists}}
  \UnaryInf$\lexists[x][!A(x)] \fCenter !A(a)$
  \RightLabel{\RightR{\lforall}}
  \UnaryInf$\lexists[x][!A(x)] \fCenter \lforall[x][!A(x)]$
  \DisplayProof\bottomAlignProof
  \qquad
  \Axiom$!A(a) \fCenter !A(a)$
  \RightLabel{*\RightR{\lforall}}
  \UnaryInf$!A(a) \fCenter \lforall[x][!A(x)]$
  \RightLabel{\LeftR{\lexists}}
  \UnaryInf$\lexists[x][!A(x)] \fCenter \lforall[x][!A(x)]$
\end{prooftree}
Padahal, $\lexists[x][!A(x)] \Sequent \lforall[x][!A(x)]$ tidak valid.
\end{explain}

\end{document}
